\SourceLeaf{273}%
\noindent If, however, we took the ten-millionth part of the Earth’s quadrant itself as the unit, on the basis of the dimensions given above, we would have
\[
\log A=6{,}8049902365
\]
\par\medskip\Needspace{4\baselineskip}\pdfbookmark[1]{Art. 6}{geodaesie1-art-6}\centerline{\textbf{6}}\nobreak\smallskip
\indent The latitude on the sphere can readily be calculated from the latitude on the ellipsoid by formula 3 if the calculation is required in only a few cases. For more extensive applications, however, it will be advantageous to use a series. The necessary formulas for developing it follow.

I denote an arbitrary latitude on the ellipsoid, or an arbitrary value of \(90^0-w\), by \(P+p\), and the corresponding latitude on the sphere, or the value of \(90^0-U\), by \(Q+q\). By \textsc{Taylor}'s theorem we have
\[
q=\frac{\mathrm{d}U}{\mathrm{d}w}\mathord{.}p-\frac{1}{2}\mathord{.}\frac{\mathrm{dd}U}{\mathrm{d}w^2}\mathord{.}pp+\frac{1}{6}\mathord{.}\frac{\mathrm{d}^3U}{\mathrm{d}w^3}\mathord{.}p^3-\frac{1}{24}\mathord{.}\frac{\mathrm{d}^4U}{\mathrm{d}w^4}\mathord{.}p^4+\text{and so on}
\]
where the derivatives are to be replaced by the particular values corresponding to \(p=0\), or to \(w=90^0-P\), \(U=90^0-Q\). Developing the derivatives successively in their general form gives
\[
\frac{\mathrm{d}U}{\mathrm{d}w}=\frac{\alpha(1-ee)\sin U}{(1-ee\cos w^2)\sin w}
\]
\[
\begin{aligned}
\frac{\mathrm{dd}U}{\mathrm{d}w^2}={}&\frac{\alpha(1-ee)\sin U}{(1-ee\cos w^2)^2\sin w^2}\\
&\quad\{\alpha(1-ee)\cos U-\cos w+ee(\cos w^3-2\cos w\sin w^2)\}
\end{aligned}
\]
\[
\begin{aligned}
\frac{\mathrm{d}^3U}{\mathrm{d}w^3}={}&\frac{\alpha(1-ee)\sin U}{(1-ee\cos w^2)^3\sin w^3}\{\alpha\alpha(1-ee)^2(\cos U^2-\sin U^2)\\
&-3\alpha(1-ee)\cos U(\cos w-ee(\cos w^3-2\cos w\sin w^2))\\
&+2\cos w^2+\sin w^2-ee(4\cos w^4-2\sin w^4)\\
&+e^4(2\cos w^6-\cos w^4\sin w^2+6\cos w^2\sin w^4)\}
\end{aligned}
\]
I have also worked out the next two derivatives, whose general forms I omit here to save space.

Substituting \(w=90^0-P\), \(U=90^0-Q\) then gives the following, provided that at the same time
\par\noindent in place of \(\alpha\sin Q\) we substitute the value \(\sin P\) (by equation 10), and
\par\noindent in place of \(\alpha\cos Q\) the value \(\frac{\cos P}{\cos\zeta\cos\eta}=\frac{\cos\theta\cos P}{\cos\varphi}\) (by equations 18, 16, 23),
\par\noindent and write \(\cos P=c\), \(\sin P=s\) for short.
